% BHU PYQ -- Transcribed & Corrected
% Paper: HINMD33: Lok Sahitya (लोक साहित्य)
% Exam: Under Graduate Semester III Examination, 2025-26 (NEP)
% Category: Multidisciplinary Course (MDC)
% Ref Code: D-III/114

\documentclass[11pt,a4paper]{article}
\usepackage[a4paper,top=2.5cm,bottom=2.5cm,left=2.5cm,right=2.5cm,headheight=14pt]{geometry}
\usepackage{amsmath,amssymb}
\usepackage[T1]{fontenc}
\usepackage[utf8]{inputenc}
\usepackage{lmodern,microtype,fancyhdr,enumitem,hyperref,lastpage,parskip}

\hypersetup{colorlinks=true,urlcolor=black,linkcolor=black,pdftitle={HINMD33: Lok Sahitya -- BHU Sem III 2025-26 (NEP MDC)}}
\pagestyle{fancy}\fancyhf{}
\fancyhead[L]{\small \textit{Banaras Hindu University}}
\fancyhead[R]{\small \textit{HINMD33: Lok Sahitya (2025-26)}}
\fancyfoot[C]{\small Page \thepage\ of \pageref{LastPage}}

\newcommand{\pts}[1]{\hfill \textit{[#1]}}

\begin{document}

\begin{center}
  {\large\bfseries BANARAS HINDU UNIVERSITY}\\[2pt]
  {\normalsize Under Graduate Semester III Examination, 2025-26 (NEP)}\\[6pt]
  \rule{0.8\linewidth}{0.5pt}\\[6pt]
  {\large\bfseries FACULTY OF ARTS / DEPARTMENT OF HINDI}\\[4pt]
  {\large\bfseries Multidisciplinary Course (MDC)}\\[4pt]
  {\large\bfseries Paper Code: HINMD33 --- लोक साहित्य (Lok Sahitya)}\\[4pt]
  {\normalsize Time: 2:30 Hours\hfill Maximum Marks: 70}
\end{center}

\rule{\linewidth}{0.4pt}\smallskip

\bigskip

\noindent \textbf{1.} किन्हीं तीन प्रश्नों के उत्तर अपने शब्दों में दीजिए । \pts{3 $\times$ 10 = 30}

\begin{enumerate}[label=(\roman*)]
  \item लोक साहित्य का अर्थ स्पष्ट करते हुए उसकी विशेषताएँ बताइए।
  \item लोक साहित्य की प्रमुख विधाओं का परिचय दीजिए।
  \item लोक कथा और लोक मिथक का अंतःसंबंध बताइए।
  \item हिन्दी की प्रमुख बोली के रचनाकारों के रूप में भिखारी ठाकुर का परिचय दीजिए।
  \item लोक साहित्य लोक संस्कृति का अभिन्न अंग है, स्पष्ट कीजिए।
\end{enumerate}

\bigskip

\noindent \textbf{2.} निम्नलिखित में से किन्हीं चार पर टिप्पणी लिखिए। \pts{4 $\times$ 6 = 24}

\begin{enumerate}[label=(\roman*)]
  \item 'लोक' का आशय
  \item लोक साहित्य और समाज
  \item लोक सुभाषित और लोक ज्ञान
  \item उद्धव शतक (जगन्नाथ दास रत्नाकर)
  \item रगई काका
  \item ईसुरी
\end{enumerate}

\bigskip

\noindent \textbf{3.} किन्हीं चार विषयों पर टिप्पणी लिखिए : \pts{4 $\times$ 4 = 16}

\begin{enumerate}[label=(\roman*)]
  \item आरसी प्रसाद सिंह का योगदान।
  \item लोकनाट्य के विविध रूप।
  \item गबर घियोर की कथावस्तु।
  \item लोककथाओं का महत्व।
  \item भारतीय संस्कृति और लोक का योगदान।
  \item भिखारी ठाकुर का योगदान।
\end{enumerate}

\bigskip
\begin{center}
  \textbf{*****}
\end{center}

\end{document}
